% Standalone XeLaTeX source; source problem: 2012ap_exam.pdf, PDF pp. 4--5.
\documentclass[11pt,a4paper]{article}
\usepackage[margin=24mm]{geometry}
\usepackage{fontspec}
\usepackage{xeCJK}
\setCJKmainfont{HaranoAjiMincho-Regular.otf}
\setCJKsansfont{HaranoAjiGothic-Medium.otf}
\usepackage{amsmath,amssymb}
\usepackage{tikz}
\usetikzlibrary{arrows.meta}
\setlength{\parindent}{1em}
\setlength{\parskip}{0.35em}
\title{物理工学専攻 院試解答集 ver2\\2012年 物理学I 第2問}
\author{政岡凜太郎}
\date{}
\begin{document}
\maketitle

陰極を $x=0$、陽極を $x=L$ とし、電位差を $V>0$ とする。電子は陰極から陽極へ進む一方、慣用電流は逆向きである。原本の装置を一次元で表すと次のようになる。
\begin{center}
\begin{tikzpicture}[>=Stealth]
  \pgfmathsetmacro{\gapLength}{5}
  \pgfmathsetmacro{\electrodeHeight}{1.8}
  \pgfmathsetmacro{\electrodeInset}{0.3}
  \pgfmathsetmacro{\electrodeMidpoint}{\electrodeHeight/2}
  \pgfmathsetmacro{\dimensionY}{-0.7*\electrodeHeight}
  \draw[very thick] (0,0)--(0,\electrodeHeight) node[above] {陰極};
  \draw[very thick] (\gapLength,0)--(\gapLength,\electrodeHeight) node[above] {陽極};
  \draw[->,thick] (\electrodeInset,\electrodeMidpoint)--(\gapLength-\electrodeInset,\electrodeMidpoint) node[midway,above] {電子 $v>0$};
  \draw[->,thick] (\gapLength-\electrodeInset,\electrodeHeight/5)--(\electrodeInset,\electrodeHeight/5) node[midway,below] {$i=-i_0$};
  \node[below=3pt,align=center] at (0,0) {$x=0$\\$\phi=0$};
  \node[below=3pt,align=center] at (\gapLength,0) {$x=L$\\$\phi=V$};
  \draw[<->] (0,\dimensionY)--(\gapLength,\dimensionY) node[midway,below] {$L$};
\end{tikzpicture}
\end{center}

\section*{[1] 電子の速度}
電子が陰極を出るときの運動エネルギーを無視できる。電荷 $-e$ の電位エネルギーは $-e\phi$ なので、陰極から $x$ までのエネルギー保存は
\begin{equation}
  \frac12 m v(x)^2-e\phi(x)=0.
\end{equation}
従って陽極向きを正とする速度は
\begin{equation}
  \boxed{v(x)=\sqrt{\frac{2e\phi(x)}{m}}}.
\end{equation}

\section*{[2] 定常電流}
電荷保存則は $\partial\rho/\partial t+\partial i/\partial x=0$ である。定常状態では $\partial\rho/\partial t=0$ だから $\partial i/\partial x=0$、すなわち $i$ は一定である。電子数密度を $n>0$、速度を $v>0$ とすると電子の電荷は $-e$ なので
\begin{equation}
  i=-env=-i_0,\qquad i_0>0.
\end{equation}

\section*{[3] 電位の微分方程式}
定常電流から $n=i_0/(ev)$ であり、これをPoisson方程式 $\epsilon_0\phi''=en$ に代入すると
\begin{align}
  \phi''(x)
  &=\frac{i_0}{\epsilon_0 v(x)}
   =\frac{i_0}{\epsilon_0}\sqrt{\frac{m}{2e}}\,\phi(x)^{-1/2}.
\end{align}
従って
\begin{equation}
  \boxed{A=\frac{i_0}{\epsilon_0}\sqrt{\frac{m}{2e}},\qquad \alpha=-\frac12.}
\end{equation}

\section*{[4] 電位分布}
$A>0$ であり、$x>0$ では $\phi$ は増加する。右辺の $\phi^{-1/2}$ は陰極端 $\phi=0$ で定義されないので、微分方程式は $x>0$ で積分し、$x\to0^+$ で $\phi,\phi'\to0$ を課して積分定数を定める。微分方程式に $\phi'$ を掛けると
\begin{equation}
  \frac12\bigl(\phi'(x)\bigr)^2
  =2A\sqrt{\phi(x)},\qquad
  \phi'(x)=2\sqrt{A}\,\phi(x)^{1/4}.
\end{equation}
この式をさらに積分して
\begin{equation}
  \frac43\phi(x)^{3/4}=2\sqrt{A}\,x
\end{equation}
を得る。よって陰極境界を満たす電位は
\begin{equation}
  \boxed{\phi(x)=\left(\frac{9A}{4}\right)^{2/3}x^{4/3}
  =\left[\frac{9i_0}{4\epsilon_0}\sqrt{\frac{m}{2e}}\right]^{2/3}x^{4/3}.}
\end{equation}
実際、この解は $\phi(0)=0$ および $\phi'(0)=0$ を満たし、$x>0$ で $\phi''>0$ となる。
ただし $\phi''(x)=\frac49(9A/4)^{2/3}x^{-2/3}$ は $x\to0^+$ で発散する。したがって微分方程式は $x>0$ で用い、陰極での電位と電場の条件は $x\to0^+$ の極限として課している。

\section*{[5] 電流電圧則}
陽極で $\phi(L)=V$ とおくと、電位分布を $i_0$ について解いて
\begin{equation}
  \boxed{i_0=\frac49\epsilon_0\sqrt{\frac{2e}{m}}\,
    \frac{V^{3/2}}{L^2}.}
\end{equation}
これは平板真空ダイオードの空間電荷制限電流密度である。$i_0$ は電圧の $3/2$ 乗に比例し、電極間隔の二乗に反比例する。符号付き慣用電流は $i=-i_0$ である。

\end{document}
